\documentclass[11pt,a4paper]{amsart}
\pdfinfoomitdate=1
\pdftrailerid{}
\pdfsuppressptexinfo=15
\usepackage[T1]{fontenc}
\usepackage{lmodern}
\usepackage[margin=27mm]{geometry}
\usepackage{microtype}
\usepackage{amsmath,amssymb,mathtools}
\usepackage{booktabs}
\usepackage{xcolor}
\usepackage[numbers,sort&compress]{natbib}
\usepackage[colorlinks=true,linkcolor=blue!55!black,citecolor=blue!55!black,urlcolor=blue!55!black]{hyperref}
\usepackage[nameinlink,noabbrev]{cleveref}
\raggedbottom

\newtheorem{theorem}{Theorem}[section]
\newtheorem{proposition}[theorem]{Proposition}
\newtheorem{corollary}[theorem]{Corollary}
\newtheorem{lemma}[theorem]{Lemma}
\theoremstyle{definition}
\newtheorem{definition}[theorem]{Definition}
\theoremstyle{remark}
\newtheorem{remark}[theorem]{Remark}

\newcommand{\Prob}{\mathbb P}
\newcommand{\Exp}{\mathbb E}
\newcommand{\N}{\mathbb N}
\newcommand{\R}{\mathbb R}
\newcommand{\one}{\mathbf 1}
\newcommand{\norm}[1]{\lVert #1\rVert}
\newcommand{\area}{C}
\newcommand{\gauss}[2]{\genfrac{[}{]}{0pt}{}{#1}{#2}_{z}}

\title{Positive Heisenberg Word-Area Cocycles}
\author{Anonymous}
\date{Internal Stage 2 draft, 29 August 2026}
\subjclass[2020]{37H15, 60F05, 20F18, 05A30}
\keywords{random matrix product, Heisenberg group, inversion statistic,
Gaussian binomial coefficient, polynomial growth, annealed pressure}

\begin{document}

\begin{abstract}
Let $X=I+E_{12}$ and $Y=I+E_{23}$, choose $X$ with probability $p$ and
$Y$ with probability $1-p$, and form the chronological left product
$M_n=A_n\cdots A_1$.  Every finite product has a three-coordinate normal
form: its central entry is the number $C_n$ of pairs in which a $Y$ occurs
before a later $X$.  Conditioning on the number of $X$ letters gives the
Gaussian-binomial inversion polynomial, while averaging the slices gives
an exact biased finite-time law.  We record
\[
 \Exp C_n=\frac{n(n-1)}2p(1-p)
\]
and a closed variance, then prove a strong law and an $n^{3/2}$ central
limit theorem with variance
$p(1-p)(3p^2-3p+1)/3$.  The same normal form yields a sharp growth
boundary: every nondegenerate Bernoulli environment has matrix-norm
exponent two in $\log n$ scale, whereas both deterministic endpoints have
exponent one.  Finally, the $n^2$-scale annealed area pressure is zero for
nonpositive tilt and equals one quarter of the positive tilt, producing a
strict gap from the typical area rate.  Gaussian-binomial inversion laws,
random-word limit theory, and random walks on Heisenberg groups are treated
as established background.  The note claims only this owner-subtracted
exact conjunction; external circulation remains on hold.
\end{abstract}

\maketitle

\section{Introduction}

Two elementary unipotent shears already retain the order of the random
letters that generate their product.  For the positive Heisenberg pair
\[
 X=I+E_{12},\qquad Y=I+E_{23},
\]
the two first-superdiagonal entries count letters.  The central entry is
more informative: it counts every $Y$ that precedes a later $X$.  Thus a
noncommutative matrix coordinate becomes an oriented area of a Bernoulli
word.

The ingredients behind this reduction have substantial prior ownership.
Gaussian polynomials enumerate inversions in fixed-content words, and their
asymptotic normality is developed by Canfield, Janson, and Zeilberger
\cite{CanfieldJansonZeilberger2011,CanfieldJansonZeilberger2012}.
For the fair binary law, the same statistic is the area under a uniformly
random north--east lattice path: Tak\'acs proved moment and limit results for
that area \cite{Takacs1986}, while Janson gave random-word inversion
interpretations, exact moments, limit theorems, and a Hoeffding-decomposition
route for uniform alphabets, including the binary case
\cite{Janson2012}.  General iid random-word subsequence methods are also
established \cite{IslakOzdemir2018}.  Random walks on
Heisenberg and higher unitriangular groups have a broad coordinate-level
limit theory \cite{DiaconisHough2021}.  The inversion law, the general
random-word CLT mechanism, the fair binary specialization, and Heisenberg
random-walk theory are not claims of this note.

The residual calculation is a finite and auditable conjunction for the
specific positive pair above.  It consists of four parts.
\begin{enumerate}
\item Every word has an exact matrix normal form, and its complete biased
      area distribution is a Bernoulli mixture of Gaussian polynomials.
\item The first two moments are closed for arbitrary $p$, and a direct
      centered-pair decomposition proves the strong law and an explicit
      $n^{3/2}$ CLT without importing a group-level limit theorem.
\item The matrix norm has polynomial exponent two for $0<p<1$ but exponent
      one at $p=0,1$.
\item Rare ordered words determine a kinked $n^2$-scale annealed pressure
      and a strict annealed--typical gap for every nonzero tilt.
\end{enumerate}

Two proof lanes are kept separate.  The first uses literal matrix
multiplication and a last-letter Gaussian-polynomial recurrence.  The
second expands the area into centered iid Bernoulli variables.  Their
intersection at the exact biased moments supplies an internal consistency
check rather than a novelty argument.  External release is \textbf{HOLD};
no absolute novelty or priority statement is made.

\section{Cocycle convention and finite-word normal form}
\label{sec:normal}

For $a,b,c\in\R$, write
\begin{equation}\label{eq:H-coordinate}
 H(a,b,c)=
 \begin{pmatrix}1&a&c\\0&1&b\\0&0&1\end{pmatrix}.
\end{equation}
The multiplication rule is
\begin{equation}\label{eq:H-law}
 H(a,b,c)H(a',b',c')
 =H(a+a',b+b',c+c'+ab').
\end{equation}
Set $X=H(1,0,0)$ and $Y=H(0,1,0)$.  Let $(A_t)_{t\geq1}$ be iid with
\begin{equation}\label{eq:environment-law}
 \Prob(A_t=X)=p,\qquad \Prob(A_t=Y)=q:=1-p,
 \qquad 0\leq p\leq1,
\end{equation}
and fix the chronological convention
\begin{equation}\label{eq:left-product}
 M_0=I,\qquad M_n=A_nA_{n-1}\cdots A_1.
\end{equation}
Thus $A_1$ acts first.

Encode $X$ by $B_t=1$ and $Y$ by $B_t=0$.  Define
\begin{equation}\label{eq:counts}
 J_n=\sum_{t=1}^n B_t,\qquad K_n=n-J_n,
\end{equation}
and the oriented word-area
\begin{equation}\label{eq:area}
 C_n=\sum_{1\leq i<j\leq n}(1-B_i)B_j.
\end{equation}
Equivalently, $C_n$ is the number of $Y$-before-$X$ pairs.  If the word is
drawn as a north--east lattice path, it is the sum of the heights of the
east steps.

\begin{theorem}[Exact normal form and slice law]
\label{thm:normal-law}
For every finite word and every $n\geq0$,
\begin{equation}\label{eq:normal-form}
 \boxed{\quad M_n=H(J_n,K_n,C_n).\quad}
\end{equation}
Moreover, for $0\leq j\leq n$,
\begin{equation}\label{eq:gaussian-slice}
 \sum_{c\geq0}
 \#\{w:J_n(w)=j,\ C_n(w)=c\}\,z^c
 =\gauss{n}{j},
\end{equation}
where the right side is the Gaussian binomial polynomial.  Consequently,
the formal probability generating polynomial is
\begin{equation}\label{eq:biased-pgf}
 \boxed{\quad
 \Exp z^{C_n}=\sum_{j=0}^n p^j q^{n-j}\gauss{n}{j}.
 \quad}
\end{equation}
The formula includes $p=0,1$, with the natural convention $0^0=1$.
\end{theorem}

\begin{proof}[Proof route I: matrix induction and last-letter recurrence]
The claim holds at time zero.  Suppose
$M_n=H(J_n,K_n,C_n)$.  Direct left multiplication gives
\begin{align}
 XM_n&=H(J_n+1,K_n,C_n+K_n),\label{eq:X-update}\\
 YM_n&=H(J_n,K_n+1,C_n).\label{eq:Y-update}
\end{align}
Appending an $X$ adds one pair for each of the $K_n$ earlier $Y$ letters;
appending a $Y$ adds no $Y$-before-$X$ pair.  These are precisely the
updates of \eqref{eq:counts}--\eqref{eq:area}, proving
\eqref{eq:normal-form}.

Let $G_{n,j}(z)$ denote the left side of \eqref{eq:gaussian-slice}.
Splitting words by their last letter gives
\begin{equation}\label{eq:gaussian-recurrence}
 G_{n,j}(z)=G_{n-1,j}(z)+z^{n-j}G_{n-1,j-1}(z),
\end{equation}
with $G_{0,0}=1$ and zero outside $0\leq j\leq n$.  This is the standard
Gaussian-binomial recurrence, so \eqref{eq:gaussian-slice} follows.
Every word in the $j$th slice has probability $p^jq^{n-j}$; summing the
slices proves \eqref{eq:biased-pgf}.
\end{proof}

The Gaussian-polynomial statement in \cref{thm:normal-law} is included to
make the cocycle calculation self-contained.  Its combinatorial content is
owned background \cite{CanfieldJansonZeilberger2011}; the fair Bernoulli
specialization is likewise direct prior territory
\cite{Takacs1986,Janson2012}.

\section{Exact biased moments and limit laws}
\label{sec:moments}

The conditional law gives one route to the exact moments.  A direct iid
expansion then proves the asymptotic statements independently of that
enumeration.

\begin{lemma}[Conditional slice moments]
\label{lem:slice-moments}
For $J_n=j$,
\begin{equation}\label{eq:conditional-moments}
 \Exp(C_n\mid J_n=j)=\frac{j(n-j)}2,
 \qquad
 \operatorname{Var}(C_n\mid J_n=j)
 =\frac{j(n-j)(n+1)}{12}.
\end{equation}
\end{lemma}

\begin{proof}
Put $k=n-j$ and $z=e^t$.  The product form of the Gaussian polynomial
gives
\[
 \frac{\gauss{n}{j}\big|_{z=e^t}}{\binom nj}
 =\prod_{r=1}^j
 \frac{(e^{(k+r)t}-1)/((k+r)t)}{(e^{rt}-1)/(rt)}.
\]
For fixed $a>0$,
\[
 \log\frac{e^{at}-1}{at}=\frac{a}{2}t+\frac{a^2}{24}t^2+O(t^3).
\]
The first two derivatives of the logarithm of the conditional moment
generating function are therefore
\[
 \frac12\sum_{r=1}^j[(k+r)-r]=\frac{jk}{2}
\]
and
\[
 \frac1{12}\sum_{r=1}^j[(k+r)^2-r^2]
 =\frac{jk(j+k+1)}{12}.
\]
These derivatives are the first two cumulants, which proves
\eqref{eq:conditional-moments}.
\end{proof}

\begin{theorem}[Biased moments, strong law, and CLT]
\label{thm:moments-clt}
For every $0\leq p\leq1$ and $n\geq0$,
\begin{equation}\label{eq:mean}
 \Exp C_n=\frac{n(n-1)}2pq
\end{equation}
and
\begin{equation}\label{eq:variance}
 \boxed{\quad
 \operatorname{Var}(C_n)
 =\frac{n(n-1)pq}{6}
 \left(6np^2-6np+2n-9p^2+9p-1\right).
 \quad}
\end{equation}
For $0<p<1$,
\begin{equation}\label{eq:slln}
 \frac{C_n}{n^2}\longrightarrow\frac{pq}{2}
 \qquad\text{almost surely},
\end{equation}
and
\begin{equation}\label{eq:clt}
 \frac{C_n-\frac{n(n-1)}2pq}{n^{3/2}}
 \Longrightarrow N(0,\sigma_p^2),
 \qquad
 \sigma_p^2=\frac{pq(3p^2-3p+1)}{3}.
\end{equation}
At $p=0,1$, $C_n=0$ deterministically.
\end{theorem}

\begin{proof}
Conditioning on $J_n\sim\operatorname{Bin}(n,p)$ and using
\cref{lem:slice-moments} proves \eqref{eq:mean}.  The law of total variance
gives
\[
 \operatorname{Var}(C_n)
 =\frac{n+1}{12}\Exp[J_n(n-J_n)]
 +\frac14\operatorname{Var}(J_n(n-J_n)).
\]
Substitution of the binomial falling moments
$\Exp[(J_n)_r]=(n)_rp^r$, for $1\leq r\leq4$, reduces the right side to
\eqref{eq:variance}.

There is also a direct derivation that does not use the conditional
Gaussian-polynomial law.  Put
$Z_{ij}=(1-B_i)B_j$ for $i<j$.  Then
$\operatorname{Var}(Z_{ij})=pq(1-pq)$.  For every triple $i<j<k$, the
three possible covariances between $Z_{ij},Z_{ik},Z_{jk}$ are
\[
 p^3q,\qquad pq^3,\qquad -p^2q^2,
\]
respectively.  Pair variables with disjoint index sets are independent.
Consequently,
\begin{equation}\label{eq:pair-covariance-variance}
 \operatorname{Var}(C_n)
 =\binom n2pq(1-pq)
  +2\binom n3pq(1-3pq),
\end{equation}
and elementary simplification gives \eqref{eq:variance}.  Thus the full
finite-time variance already has two independent derivations.

For an independent route to the limits, put $\eta_k=B_k-p$.  Expanding
each summand in \eqref{eq:area} yields the exact centered identity
\begin{equation}\label{eq:centered-decomposition}
 C_n-\frac{n(n-1)}2pq
 =L_n+R_n,
\end{equation}
where
\begin{align}
 L_n&=\sum_{k=1}^n[k-1-p(n-1)]\eta_k,
 \label{eq:linear-part}\\
 R_n&=-\sum_{1\leq i<j\leq n}\eta_i\eta_j
 =-\frac12\left[
 \left(\sum_{k=1}^n\eta_k\right)^2-\sum_{k=1}^n\eta_k^2
 \right].\label{eq:quadratic-part}
\end{align}

The ordinary strong law gives
$S_n:=\sum_{k=1}^n\eta_k=o(n)$ almost surely.  Summation by parts gives
\[
 \sum_{k=1}^n(k-1)\eta_k
 =(n-1)S_n-\sum_{m=1}^{n-1}S_m=o(n^2).
\]
Equations \eqref{eq:linear-part}--\eqref{eq:quadratic-part} then imply
$L_n=o(n^2)$ and $R_n=o(n^2)$ almost surely.  Combining this with
\eqref{eq:centered-decomposition} proves \eqref{eq:slln}.

For the CLT, independence gives
\begin{align}
 \operatorname{Var}(L_n)
 &=pq\sum_{k=1}^n[k-1-p(n-1)]^2\notag\\
 &=\frac{pq\,n(n-1)}6
 \left(2n-1-6p(n-1)+6p^2(n-1)\right).
 \label{eq:linear-variance}
\end{align}
After division by $n^3$, this converges to $\sigma_p^2>0$.  Every summand
in \eqref{eq:linear-part} has magnitude at most $n$, whereas
$\sqrt{\operatorname{Var}(L_n)}$ has order $n^{3/2}$.  The Lindeberg
condition is therefore eventually vacuous, and the triangular-array
Lindeberg--Feller theorem gives
$L_n/n^{3/2}\Longrightarrow N(0,\sigma_p^2)$.
Finally,
\[
 \Exp|R_n|
 \leq\frac12\left(
 \Exp S_n^2+\sum_{k=1}^n\Exp\eta_k^2\right)=npq.
\]
Hence $R_n/n^{3/2}\to0$ in probability.  Slutsky's theorem proves
\eqref{eq:clt}.  The endpoint statement follows directly from
\eqref{eq:area}.
\end{proof}

At $p=1/2$, \eqref{eq:variance} becomes
$n(n-1)(2n+5)/96$, while $\sigma_{1/2}^2=1/48$.

\section{Polynomial matrix growth}
\label{sec:growth}

All norms on the nine-dimensional space of $3\times3$ real matrices are
equivalent.  The positive normal form therefore turns area growth into
matrix-norm growth without cancellation.

\begin{theorem}[Sharp exponent boundary]
\label{thm:norm-boundary}
For every fixed matrix norm,
\begin{equation}\label{eq:norm-exponent}
 \lim_{n\to\infty}\frac{\log\norm{M_n}}{\log n}
 =\begin{cases}
 2,&0<p<1,\\
 1,&p=0\text{ or }p=1,
 \end{cases}
 \qquad\text{almost surely}.
\end{equation}
\end{theorem}

\begin{proof}
It suffices to use the Frobenius norm, for which
\begin{equation}\label{eq:Frobenius}
 \norm{M_n}_{\mathrm F}^2=3+J_n^2+K_n^2+C_n^2.
\end{equation}
If $0<p<1$, \eqref{eq:slln} gives
$C_n/n^2\to pq/2>0$ almost surely.  Equation \eqref{eq:Frobenius} then
implies $\norm{M_n}_{\mathrm F}=n^{2+o(1)}$.  At $p=1$,
$M_n=I+nE_{12}$; at $p=0$, $M_n=I+nE_{23}$.  In either endpoint,
$\norm{M_n}_{\mathrm F}^2=n^2+3$.  Norm equivalence changes logarithms by a
bounded amount and therefore preserves \eqref{eq:norm-exponent}.
\end{proof}

This boundary is polynomial rather than Lyapunov: all eigenvalues of every
$M_n$ equal one, and $n^{-1}\log\norm{M_n}\to0$ in every regime.

\section{Quadratic area pressure}
\label{sec:pressure}

For $\theta\in\R$, define the annealed area pressure and its pathwise
counterpart by
\begin{align}
 \mathcal P_p(\theta)
 &=\lim_{n\to\infty}\frac1{n^2}
   \log\Exp e^{\theta C_n},\label{eq:annealed-pressure}\\
 \mathcal Q_p(\theta)
 &=\lim_{n\to\infty}\frac{\theta C_n}{n^2},
 \label{eq:typical-pressure}
\end{align}
whenever the limits exist.  The $n^2$ normalization matches the central
coordinate, rather than the usual exponential norm scale.

\begin{proposition}[Extremal probabilities]
\label{prop:extrema}
For every $n\geq0$,
\begin{equation}\label{eq:max-area}
 \max_w C_n(w)=\left\lfloor\frac{n^2}{4}\right\rfloor.
\end{equation}
For $0<p<1$,
\begin{equation}\label{eq:max-probability}
 \Prob\left(C_n=\left\lfloor\frac{n^2}{4}\right\rfloor\right)
 =(pq)^{\lfloor n/2\rfloor}.
\end{equation}
The zero-area probability is
\begin{equation}\label{eq:zero-probability}
 \Prob(C_n=0)=\sum_{j=0}^n p^jq^{n-j}
 =\begin{cases}
 \dfrac{p^{n+1}-q^{n+1}}{p-q},&p\neq q,\\[6pt]
 (n+1)2^{-n},&p=q=\frac12.
 \end{cases}
\end{equation}
\end{proposition}

\begin{proof}
With $j$ letters $X$ and $n-j$ letters $Y$, at most $j(n-j)$ ordered
pairs contribute.  Equality requires the word $Y^{n-j}X^j$.  Maximizing
$j(n-j)$ over integers gives \eqref{eq:max-area}.  There is one balanced
maximizer when $n$ is even and two when $n$ is odd; their probabilities
sum to \eqref{eq:max-probability}.  Area zero means that no $Y$ precedes an
$X$, so the word is uniquely $X^jY^{n-j}$ for some $0\leq j\leq n$.
Summing their probabilities gives \eqref{eq:zero-probability}.
\end{proof}

\begin{theorem}[Pressure kink and strict gap]
\label{thm:pressure}
The limit \eqref{eq:annealed-pressure} exists for every $p\in[0,1]$ and
$\theta\in\R$.  It is
\begin{equation}\label{eq:pressure-formula}
 \boxed{\quad
 \mathcal P_p(\theta)=
 \begin{cases}
 \theta/4,&0<p<1\text{ and }\theta>0,\\
 0,&0<p<1\text{ and }\theta\leq0,\\
 0,&p\in\{0,1\}.
 \end{cases}\quad}
\end{equation}
For $0<p<1$, the pathwise limit is
\begin{equation}\label{eq:typical-formula}
 \mathcal Q_p(\theta)=\frac{\theta pq}{2}
 \qquad\text{almost surely}.
\end{equation}
Thus $\mathcal P_p(\theta)>\mathcal Q_p(\theta)$ for every
$\theta\neq0$ and $0<p<1$.
\end{theorem}

\begin{proof}
Suppose first that $0<p<1$ and $\theta>0$.  By
\cref{prop:extrema},
\[
 (pq)^{\lfloor n/2\rfloor}
 e^{\theta\lfloor n^2/4\rfloor}
 \leq \Exp e^{\theta C_n}
 \leq e^{\theta\lfloor n^2/4\rfloor}.
\]
Taking logarithms and dividing by $n^2$ proves the first line of
\eqref{eq:pressure-formula}.  If $\theta<0$, then
\[
 \Prob(C_n=0)\leq\Exp e^{\theta C_n}\leq1.
\]
The probability in \eqref{eq:zero-probability} is bounded below by
$\max\{p,q\}^n$, so both logarithmic bounds have limit zero on the
$n^2$ scale.  The case $\theta=0$ is immediate.  At $p=0,1$, $C_n=0$.

Equation \eqref{eq:typical-formula} is \eqref{eq:slln} multiplied by
$\theta$.  For positive $\theta$,
$1/4-pq/2\geq1/8$; for negative $\theta$, the annealed value is zero while
the typical value is negative.  This proves strictness.
\end{proof}

\section{Proof lanes, controls, and owner boundary}
\label{sec:scope}

The two proof lanes use different primitive data.
\begin{itemize}
\item \textbf{Finite-word lane.}  Raw matrix multiplication gives
      \eqref{eq:normal-form}; a last-letter recurrence gives the complete
      conditional area polynomial.  Differentiating its product form gives
      the conditional moments, and binomial mixing gives the biased moments.
\item \textbf{Centered-pair lane.}  The iid identity
      \eqref{eq:centered-decomposition} does not use Gaussian polynomials.
      Shared-index covariance counting first recovers the complete variance
      \eqref{eq:variance}; the centered identity then proves the strong law
      and CLT.  This lane meets the first at an exact finite-time formula,
      not only at its leading coefficient.
\end{itemize}
The norm theorem uses positivity and finite-dimensional norm equivalence.
The pressure theorem instead uses exact extremizers and their probabilities;
neither conclusion is extrapolated from computation.

The standard-library verifier mirrors these lanes.  It multiplies all
$131{,}071$ binary words through length $16$, independently scans their
letter and area statistics, and compares every conditional histogram with
the Gaussian recurrence.  Separate exact-rational lanes check biased moments
and the independent pair-covariance formula through time $32$, the centered
decomposition word by word through time $12$, endpoint norms, extremal
probabilities, and exponential-moment bounds.
The stored output records the exact assertion count.  These controls are
finite falsification tests, not proofs of the asymptotic statements.

The owner boundary is part of the result's scope.  Gaussian-binomial
enumeration and fixed-content inversion asymptotics are subtracted against
\cite{CanfieldJansonZeilberger2011,CanfieldJansonZeilberger2012}; fair-word
inversion moments and limit laws are subtracted directly against
\cite{Takacs1986,Janson2012}; and general iid random-word subsequence methods
are subtracted against \cite{IslakOzdemir2018}.  General Heisenberg and
unitriangular random-walk limit theory is subtracted against
\cite{DiaconisHough2021}.  The internally studied residual package is only
the arbitrary-bias conjunction for this particular matrix cocycle, the
polynomial exponent boundary, and the quadratic pressure kink for the two
displayed positive generators.  This is a scope description, not a novelty
claim.  Search absence is not evidence, and external circulation remains
\textbf{HOLD} pending specialist direct-owner review.

Internally, P70 studies weighted shifts on finite Heisenberg quotients via
convolution nullities; P93 is a symbolic push--pop stack cocycle controlled
by a reflected random walk; P99 studies one deterministic shear acting on
finite-index sublattices; and P104 studies exponentially contracting
monomial random matrices through occupation and singular-value pressure.
The phase spaces and observables here are different, but the shared cocycle,
Heisenberg, and random-matrix vocabulary remains disclosed.

\section{Conclusion}

\enlargethispage{3\baselineskip}

The positive Heisenberg product stores Bernoulli word order in one central
entry.  Combining its exact finite-word law with an independent centered
pair expansion yields the complete biased moments, a strong law, and an
explicit CLT.  Positivity then exposes a polynomial norm boundary, while
rare ordered words force a kinked quadratic pressure and a strict
annealed--typical gap.  These conclusions define the paper's full scope;
all broader inversion, random-word, and Heisenberg-walk theory remains
owned background.  External release is on hold.

\bibliographystyle{plainnat}
\bibliography{references}

\end{document}
